\small
\setlength{\extrarowheight}{2pt}
\renewcommand{\arraystretch}{1.2}

\begin{longtable}{
    >{\raggedright\arraybackslash}p{2.8cm}
    >{\raggedright\arraybackslash}p{2.8cm}
    >{\raggedright\arraybackslash}p{6.2cm}
    >{\raggedright\arraybackslash}p{2.2cm}
}

\caption{Ferramentas Open Source Identificadas na Literatura}
\label{tab:ferramentas_open_source} \\

\toprule
\textbf{Ferramenta} &
\textbf{Contexto de Uso} &
\textbf{Resumo} &
\textbf{Citações} \\
\midrule
\endfirsthead

\multicolumn{4}{c}{\textit{Continuação da Tabela \ref{tab:ferramentas_open_source}}} \\
\toprule
\textbf{Ferramenta} &
\textbf{Contexto de Uso} &
\textbf{Resumo} &
\textbf{Citações} \\
\midrule
\endhead

\bottomrule
\endlastfoot

\rowcolor[gray]{0.96}
Apache Kafka &
Ingestão de dados &
Processamento de fluxos de dados em tempo real, com alta escalabilidade e tolerância a falhas. &
\cite{murarka2024, mbata2024survey} \\

Apache Flume &
Ingestão de dados &
Coleta e transferência de logs e eventos para ambientes Big Data. &
\cite{murarka2024} \\

\rowcolor[gray]{0.96}
Apache NiFi &
Ingestão e ETL &
Automação de fluxos de dados, ingestão, transformação e integração entre sistemas. &
\cite{murarka2024, mbata2024survey} \\

RabbitMQ &
Fila de mensagens &
Troca assíncrona de mensagens entre aplicações e serviços distribuídos. &
\cite{septian2019} \\

\rowcolor[gray]{0.96}
Apache Spark &
Processamento de dados &
Processamento distribuído em memória para cargas batch e streaming. &
\cite{murarka2024, septian2019, mbata2024survey} \\

Apache Airflow &
Orquestração &
Gerenciamento, agendamento e monitoramento de pipelines de dados e ML. &
\cite{mbata2024survey} \\

\rowcolor[gray]{0.96}
HDFS &
Armazenamento &
Base para Data Lakes, oferecendo armazenamento distribuído e tolerante a falhas. &
\cite{murarka2024, septian2019} \\

Apache Zeppelin &
Visualização e BI &
Criação de dashboards, análises interativas e exploração de dados. &
\cite{septian2019} \\

\rowcolor[gray]{0.96}
Metabase &
Visualização e BI &
Ferramenta de Business Intelligence para consultas, relatórios e dashboards. &
\cite{septian2019} \\

\end{longtable}